\documentclass[11pt,a4paper]{article}
\usepackage[T1]{fontenc}
\usepackage[utf8]{inputenc}
\usepackage{lmodern}
\usepackage{amsmath,amssymb,amsthm,mathtools}
\usepackage[margin=27mm,headheight=14pt]{geometry}
\usepackage{microtype}
\usepackage{enumitem}
\usepackage{xcolor}
\usepackage{fancyhdr}
\usepackage{booktabs}
\usepackage{xurl}
\usepackage{hyperref}
\usepackage{bookmark}
\usepackage[capitalise,nameinlink,noabbrev]{cleveref}
\definecolor{linkblue}{RGB}{22,68,105}
\hypersetup{colorlinks=true,linkcolor=linkblue,citecolor=linkblue,urlcolor=linkblue,
  pdftitle={The radius-two supported norm of the modular group acting on prime projective lines},
  pdfauthor={Yitzchak Shmalo},
  pdfsubject={Schreier graphs, modular group, projective line, exact certificates}}
\urlstyle{same}
\setlength{\parindent}{1.2em}
\setlength{\parskip}{0.2em}
\setlength{\emergencystretch}{1.5em}
\setlist[enumerate]{itemsep=0.35em,topsep=0.45em}
\numberwithin{equation}{section}
\newtheorem{theorem}{Theorem}[section]
\newtheorem{lemma}[theorem]{Lemma}
\newtheorem{proposition}[theorem]{Proposition}
\newtheorem{corollary}[theorem]{Corollary}
\theoremstyle{remark}
\newtheorem{remark}[theorem]{Remark}
\newcommand{\Z}{\mathbb Z}
\newcommand{\Q}{\mathbb Q}
\newcommand{\F}{\mathbb F}
\newcommand{\PP}{\mathbb P}
\newcommand{\PSL}{\mathrm{PSL}}
\newcommand{\SL}{\mathrm{SL}}
\pagestyle{fancy}
\fancyhf{}
\fancyhead[L]{\small Radius-two supported norm on prime projective lines}
\fancyhead[R]{\small 25 November 2026}
\fancyfoot[C]{\thepage}
\renewcommand{\headrulewidth}{0.4pt}

\title{The radius-two supported norm of the modular group acting on prime projective lines}
\author{Yitzchak Shmalo\\[0.4em]
\small Einstein Institute of Mathematics, The Hebrew University of Jerusalem\\
\small Givat Ram, Jerusalem, Israel}
\date{25 November 2026}

\begin{document}
\maketitle
\thispagestyle{fancy}

\begin{abstract}
Let $A_p$ be the sum of the permutation operators of the generators $S$, $R$, $R^{-1}$ of $\PSL_2(\Z)$ on $\ell^2(\PP^1(\F_p))$. We determine, for every prime $p$, the largest norm of $A_p$ on vectors supported in one ball of radius two of the Schreier graph. It is $3$ for $p\le7$, $\sqrt\beta=2.90782\ldots$ for $p=13$, and $\sqrt\alpha=2.890889\ldots$ for every other prime, where $\alpha$ and $\beta$ are the largest roots of $t^5-19t^4+124t^3-331t^2+325t-60$ and $t^5-18t^4+107t^3-242t^2+169t-24$; for the primes in the last case the maximum is attained only at the centre $1$. Outside eleven explicit primes the problem reduces to $23$ integer models built from words of length at most six, and the eleven are settled by exact computation. The result concerns a fixed radius and gives no bound on the global norm off the constants.
\end{abstract}

\noindent\textit{Keywords:} modular group, Schreier graph, projective line over a finite field, operator norm, centralizers, exact certificates.\\
\textit{2020 Mathematics Subject Classification:} 05C50, 20H05, 05C25, 47A30.

\section{Introduction}\label{sec:intro}

The modular group $\Gamma=\PSL_2(\Z)$ is the free product $C_2*C_3$, generated by the images of
\[
S=\begin{pmatrix}0&-1\\1&0\end{pmatrix},\qquad R=\begin{pmatrix}0&-1\\1&1\end{pmatrix},\qquad R^{-1}=\begin{pmatrix}1&1\\-1&0\end{pmatrix}.
\]
For a prime $p$ it acts on the projective line $X_p=\PP^1(\F_p)$, a set of $p+1$ points, by M\"obius transformations. Let $\pi_p(g)\delta_x=\delta_{gx}$ be the permutation representation on $\ell^2(X_p)$ and put
\[
A_p=\pi_p(S)+\pi_p(R)+\pi_p(R^{-1}).
\]
This is the adjacency operator of the Schreier graph of $X_p$ with respect to $\{S,R,R^{-1}\}$, counted with loops and multiple edges. Its norm on all of $\ell^2(X_p)$ is $3$, attained by the constants. The interesting quantity is its norm on the orthogonal complement of the constants. Van Handel~\cite[\S6.4]{vanHandel} poses the deterministic question whether the restriction of $\pi_p$ to the constants-orthogonal space converges strongly to the regular representation of $\Gamma$ as $p\to\infty$, a convergence conjectured by Buck~\cite{Buck} for diffusion operators; for $A_p$ that would mean the norm off the constants tends to the regular norm
\[
\|\lambda(S)+\lambda(R)+\lambda(R^{-1})\|=\tfrac12\bigl(1+\sqrt{13+8\sqrt2}\bigr)=2.9654\ldots,
\]
which goes back to the work of Cartwright and Soardi~\cite{CartwrightSoardi} on free products and is recorded for $C_2*C_3$ in~\cite{BarbieriJezernik}. Barbieri and Jezernik~\cite{BarbieriJezernik} prove strong convergence for uniformly random permutation representations of free products of finite groups; that is a different model from the prescribed sequence $X_p$. The numerical evidence for the convergence recorded in~\cite[\S6.4]{vanHandel} includes the spectral computations of Lafferty and Rockmore for the special linear group $\SL_2$ over finite fields~\cite{LaffertyRockmore92,LaffertyRockmore93}. Rivin and Sardari~\cite{RivinSardari} study the spectra of random Cayley graphs of $\SL_2(\Z/p\Z)$, and Golubev and Kamber~\cite{GolubevKamber} prove cutoff for Schreier graphs of $\SL_2(\F_t)$ acting on the projective line with random generators, through the Sarnak--Xue density of eigenvalues.

We compute a local relative of the global norm exactly. For $x\in X_p$ let $B_2(x)$ be the closed ball of radius two about $x$ in the Schreier graph, and let $P_{2,x}$ be the coordinate projection of $\ell^2(X_p)$ onto functions supported in $B_2(x)$. Define
\[
N_2(p)=\max_{x\in X_p}\|A_pP_{2,x}\|.
\]
The output of $A_pP_{2,x}$ is unrestricted, so this is not the norm of the induced adjacency matrix of the ball, and the input is not required to be orthogonal to the constants.

\begin{theorem}\label{thm:main}
Let $\alpha$ and $\beta$ be the largest real roots of
\[
Q(t)=t^5-19t^4+124t^3-331t^2+325t-60,\qquad Q_{13}(t)=t^5-18t^4+107t^3-242t^2+169t-24 .
\]
For every prime $p$,
\[
N_2(p)=\begin{cases}3,& p\in\{2,3,5,7\},\\ \sqrt\beta,& p=13,\\ \sqrt\alpha,& \text{otherwise.}\end{cases}
\]
For $p\notin\{2,3,5,7,13\}$ the maximum is attained only at $x=1$, and for $p=13$ exactly at $x=2$ and $x=7$. Moreover
\begin{gather*}
8.3572446843<\alpha<8.3572446844,\qquad 2.890889946779<\sqrt\alpha<2.890889946780,\\
8.4554193488<\beta<8.4554193489,\qquad \sqrt\beta=2.90782\ldots .
\end{gather*}
\end{theorem}

What is new here is the exact value of $N_2(p)$ at every prime, the identification of the maximizing centre, and the complete list of primes at which the value differs from $\sqrt\alpha$. The lemma that a vanishing additive commutator forces commutation (\cref{lem:commute}) and the centralizer argument of \cref{sec:models} are elementary; the common-eigenline criterion behind the first is classical (see Florentino~\cite[Proposition~2.13]{Florentino}, who attributes it to Friedland~\cite{Friedland}). The constant $\sqrt\alpha$ is below the regular norm because the support constraint is fixed; the theorem gives no bound for globally supported vectors orthogonal to the constants, and the deterministic question of~\cite{vanHandel} is not addressed.

The proof has three parts. \Cref{sec:short} shows that two short words can share a fixed point modulo $p$ without commuting only for $p$ in an explicit set $E$ of eleven primes. \Cref{sec:models} shows that outside $E$ the matrix of $A_pP_{2,x}$ depends only on which commutation class of short words fixes $x$, which leaves $23$ integer models. \Cref{sec:proof} compares the models exactly and treats the primes in $E$ by direct computation. All finite computations use integer and rational arithmetic only; the programs are described in \cref{sec:computations}.

\section{Short words and the exceptional primes}\label{sec:short}

Let $B_L\subset\Gamma$ be the set of elements of word length at most $L$ in $S,R,R^{-1}$. We represent each element by the integer matrix of determinant one whose first nonzero entry is positive. Multiplying out gives
\[
|B_0|,\ldots,|B_6|=1,4,8,14,22,34,50 .
\]
The Schreier graph has an edge from $y$ to $sy$ for each $s\in\{S,R,R^{-1}\}$; since the generating set is closed under inverses, $B_r(x)=\{gx:g\in B_r\}$.

For $G,H\in\SL_2(\Z)$ put $\Delta(G,H)=\det(GH-HG)$. Replacing $G$ or $H$ by its negative does not change $\Delta$, so $\Delta$ is defined on $\Gamma$.

\begin{lemma}\label{lem:divides}
If $G$ and $H$ fix the same point of $X_p$, then $p\mid\Delta(G,H)$.
\end{lemma}

\begin{proof}
A common fixed point is a line in $\F_p^2$ spanned by a nonzero $v$ with $Gv=\lambda v$ and $Hv=\mu v$ modulo $p$. Then $(GH-HG)v\equiv0$, so $GH-HG$ is singular modulo $p$.
\end{proof}

\begin{lemma}\label{lem:commute}
If $G,H\in\SL_2(\Z)$ and $\Delta(G,H)=0$, then $GH=HG$.
\end{lemma}

\begin{proof}
Over an algebraically closed field, $\Delta(G,H)=0$ is equivalent to $G$ and $H$ having a common eigenline (Friedland; see~\cite[Proposition~2.13]{Florentino}). We give the short argument for the direction we need and then use integrality. If $G$ is scalar there is nothing to prove, so assume it is not. Over $\mathbb C$, write $G$ in Jordan form. If $G=\mathrm{diag}(\lambda,\mu)$ with $\lambda\ne\mu$, then $GH-HG$ has zero diagonal and off-diagonal entries $(\lambda-\mu)h_{12}$ and $(\mu-\lambda)h_{21}$, so $\Delta=(\lambda-\mu)^2h_{12}h_{21}$, and $\Delta=0$ means that one eigenline of $G$ is $H$-invariant. If $G=\left(\begin{smallmatrix}\lambda&1\\0&\lambda\end{smallmatrix}\right)$ is a nontrivial Jordan block, then
\[
GH-HG=\begin{pmatrix}h_{21}&h_{22}-h_{11}\\0&-h_{21}\end{pmatrix},\qquad \Delta=-h_{21}^2,
\]
and again $\Delta=0$ means that the eigenline of $G$ is $H$-invariant.

Suppose $G$ has distinct eigenvalues. Its characteristic polynomial $t^2-(\operatorname{tr}G)t+1$ is irreducible over $\Q$, since a rational root would be an integer dividing $1$, and $\pm1$ would be a double root. The eigenlines of $G$ are therefore conjugate over the quadratic field of $G$, and since $H$ is rational, an $H$-invariant eigenline has an $H$-invariant conjugate. Both eigenlines of $G$ are $H$-invariant, $H$ is diagonal in the eigenbasis of $G$, and the two commute.

Suppose $G$ has a repeated eigenvalue. It is then $\pm1$, and its unique eigenline is rational. The eigenvalue of $H$ on that line is rational and a root of the monic integer polynomial $t^2-(\operatorname{tr}H)t+1$, hence $\pm1$, and so it is the repeated eigenvalue of $H$. In a rational basis whose first vector spans the line, $G$ and $H$ are upper triangular with constant diagonals; such matrices commute.
\end{proof}

Integrality and determinant one are both used: the rational pair $\mathrm{diag}(2,1/2)$, $\left(\begin{smallmatrix}1&1\\0&1\end{smallmatrix}\right)$, and the integral pair $\left(\begin{smallmatrix}1&0\\0&-1\end{smallmatrix}\right)$, $\left(\begin{smallmatrix}1&1\\0&-1\end{smallmatrix}\right)$ of determinant $-1$, have singular commutators and do not commute.

Two elements of $\Gamma$ commute exactly when their representatives in $\SL_2(\Z)$ commute, because $GH=-HG$ is impossible in $\SL_2(\Z)$. Indeed, taking traces in $H=-G^{-1}HG$ and $G=-H^{-1}GH$ would give $\operatorname{tr}G=\operatorname{tr}H=0$, hence $G^2=H^2=-I$, and $i\mapsto G$, $j\mapsto H$ would define an algebra homomorphism from Hamilton's quaternions $\mathbb H$ to $M_2(\mathbb R)$. Since $\mathbb H$ is a division algebra, this map is injective, and by dimension count it is onto; this is impossible, because $M_2(\mathbb R)$ has zero divisors and $\mathbb H$ has none.

Let $E$ be the set of primes $p$ for which some $x\in X_p$ is fixed by two elements of $B_6\setminus\{1\}$ that do not commute.

\begin{proposition}\label{prop:E}
$E=\{2,3,5,7,11,13,37,41,61,73,97\}$.
\end{proposition}

\begin{proof}
Among the $1{,}176$ unordered pairs of distinct elements of $B_6\setminus\{1\}$, exactly $50$ have $\Delta=0$, and the largest value of $|\Delta|$ is $100$, attained for instance by $\left(\begin{smallmatrix}1&-2\\1&-1\end{smallmatrix}\right)$ and $\left(\begin{smallmatrix}1&2\\-2&-3\end{smallmatrix}\right)$. If $p\in E$, with noncommuting $g,h\in B_6\setminus\{1\}$ fixing a point, then $\Delta(g,h)\ne0$ by \cref{lem:commute} and $p\mid\Delta(g,h)$ by \cref{lem:divides}, so $p\le 97$. A direct check of every point of $X_p$ for every prime $p\le 101$ gives the stated set. The prime divisors of the nonzero values of $\Delta$ are exactly the elements of $E$.
\end{proof}

We also need that short words stay nonscalar after reduction. For $g=\left(\begin{smallmatrix}a&b\\c&d\end{smallmatrix}\right)$ the reduction of $g$ modulo $p$ is scalar exactly when $p$ divides $\gcd(b,c,a-d)$. For the $49$ elements of $B_6\setminus\{1\}$ the only primes dividing one of these gcds are $2$ and $3$, so every element of $B_6\setminus\{1\}$ is nonscalar modulo every $p\notin E$.

\section{The local models}\label{sec:models}

For an integer matrix $M=\left(\begin{smallmatrix}a&b\\c&d\end{smallmatrix}\right)$ put $v(M)=(b,c,a-d)$. A direct expansion shows that the entries of $MN-NM$ are $\pm$ the $2\times2$ minors of the matrix with rows $v(M)$ and $v(N)$:
\[
MN-NM=\begin{pmatrix} bc'-b'c & b'(a-d)-b(a'-d')\\ c(a'-d')-c'(a-d) & cb'-bc'\end{pmatrix}.
\]
Hence two nonscalar matrices commute exactly when their vectors $v$ are proportional. Sending a nonscalar matrix to the primitive integer vector proportional to $v(M)$ with positive first nonzero coordinate therefore partitions $B_6\setminus\{1\}$ into commutation classes, each a set of pairwise commuting elements, no two classes commuting. There are $22$ classes.

For $p$ prime and $x\in X_p$ let $F_x=\{g\in B_6\setminus\{1\}:gx=x\}$.

\begin{lemma}\label{lem:class}
If $p\notin E$, then $F_x$ is empty or equal to one commutation class.
\end{lemma}

\begin{proof}
Elements of $F_x$ commute pairwise because $p\notin E$, so a nonempty $F_x$ lies in one class $C$. Conversely let $g\in F_x$ and $h\in C$. The reductions $\bar g,\bar h$ commute and are nonscalar since $p>3$. The centralizer of a nonscalar $2\times2$ matrix $T$ over a field is $\{\lambda I+\mu T\}$, so $\bar h=\lambda I+\mu\bar g$ with $\mu\ne0$, and $\bar h$ and $\bar g$ have the same invariant lines. Points of $X_p$ fixed by $g$ are its invariant lines, so $h$ fixes $x$.
\end{proof}

For $g,h\in B_3$ we have $gx=hx$ exactly when $h^{-1}g$ fixes $x$, and $h^{-1}g\in B_6$. So, for $p\notin E$, the equivalence relation $g\sim h\iff h^{-1}g\in\{1\}\cup F_x$ on $B_3$ describes all coincidences among the points $gx$, $g\in B_3$. The ball $B_2(x)$ is the set of classes meeting $B_2$, and the vertices $sy$ with $y\in B_2(x)$ are classes of $B_3$. Choose one element of $B_2$ in each class meeting $B_2$. The matrix of $A_pP_{2,x}$, restricted to the nonzero rows, has rows indexed by the classes of $B_3$ that are hit and columns indexed by the chosen representatives, and its entry at a row $y$ and a representative $g$ is the number of $s\in\{S,R,R^{-1}\}$ with $sg\in y$. Equal vertices give one coordinate, and loops and multiple edges are counted.

For each commutation class $C$, and for $C=\varnothing$, let $M_C$ be the integer matrix obtained in this way with $F_x$ replaced by $C$, and let $G_C=M_C^{\mathsf T}M_C$. Then $\|A_pP_{2,x}\|^2$ is the largest eigenvalue of $G_{F_x}$, the matrices $G_C$ do not depend on $p$, and there are $23$ of them, of size at most $8$. Nothing requires every model to occur at every prime.

\section{Proof of the theorem}\label{sec:proof}

We first show that one model dominates the others, then deduce the theorem for primes outside $E$, and finally treat the eleven primes of $E$.

\subsection{The maximal model}

Let $C_*$ be the class containing
\[
g_1=\begin{pmatrix}0&1\\-1&2\end{pmatrix},\qquad g_2=\begin{pmatrix}2&-1\\1&0\end{pmatrix},
\]
both of word length five; these are its only elements. In a suitable ordering of the eight input coordinates,
\[
G_{C_*}=\begin{pmatrix}
3&1&0&0&1&1&1&0\\
1&3&2&2&0&0&1&0\\
0&2&5&1&1&0&0&0\\
0&2&1&5&1&0&0&0\\
1&0&1&1&3&1&1&1\\
1&0&0&0&1&3&0&0\\
1&1&0&0&1&0&3&1\\
0&0&0&0&1&0&1&3
\end{pmatrix},
\]
with characteristic polynomial
\[
t^8-28t^7+320t^6-1942t^5+6784t^4-13740t^3+15285t^2-8000t+1200=(t-4)(t^2-5t+5)\,Q(t).
\]
The roots of $t^2-5t+5$ are below $4$. All leading principal minors of $u^2I-G_{C_*}$ are positive at $u=2.890889946780$, so every eigenvalue is below $u^2$, and $\det(l^2I-G_{C_*})<0$ at $l=2.890889946779$, so some eigenvalue exceeds $l^2>4$. That eigenvalue is a root of $Q$; since $G_{C_*}$ is symmetric, all roots of $Q$ are eigenvalues, so the largest eigenvalue of $G_{C_*}$ is $\alpha$ and $l<\sqrt\alpha<u$. The same two tests at $8.3572446844$ and $8.3572446843$ give the bracket for $\alpha$.

For each of the other $22$ models all leading principal minors of $83I-10G_C$ are positive (the smallest of these integers is $33$), so by Sylvester's criterion the largest eigenvalue of $G_C$ is below $8.3<\alpha$.

\subsection{Primes outside \texorpdfstring{$E$}{E}}

Let $p\notin E$. By \cref{lem:class} and \cref{sec:models}, $N_2(p)^2$ is the largest eigenvalue of $G_C$ over the classes $C$ (or $\varnothing$) that occur as some $F_x$. Both $g_1$ and $g_2$ fix $x=1$: $g_1(1)=1/(2-1)$ and $g_2(1)=(2-1)/1$. Hence $F_1=C_*$ and $N_2(p)=\sqrt\alpha$. If the maximum is attained at $x$, then $F_x=C_*$ since every other model is smaller, so $g_1$ fixes $x$. The point $\infty$ is not fixed, since the lower-left entry of $g_1$ is $-1$, and for finite $x$ the equation $x=1/(2-x)$ is $(x-1)^2=0$. So $x=1$ is the only maximizing centre.

\subsection{The primes in \texorpdfstring{$E$}{E}}

For $p\in E$ we build the integer matrix of $A_pP_{2,x}$ directly over $\F_p$ at every $x\in X_p$ and compute the characteristic polynomial of its Gram matrix exactly.

For $p=2,3,5,7$ some ball $B_2(x)$ is all of $X_p$: every $x$ for $p=2,3$, the points $1$ and $4$ for $p=5$, and the point $1$ for $p=7$. Then $P_{2,x}$ is the identity and the constant vector gives $\|A_pP_{2,x}\|=3$; on the other hand $\|A_p\|\le3$ since each $\pi_p(g)$ is unitary.

For $p=11,37,41,61,73,97$ the polynomial $(t-4)(t^2-5t+5)Q(t)$ occurs at $x=1$ and nowhere else, and at every other centre all leading principal minors of $8.3572446843\,I-G$ are positive. So these primes satisfy the generic formula with the unique centre $x=1$.

For $p=13$ the polynomial $(t-3)(t-2)Q_{13}(t)$ occurs exactly at $x=2$ and $x=7$. At every other centre all leading principal minors of $\frac{42}{5}I-G$ are positive, so the squared norm there is below $8.4$. At $x=2$, all leading minors of $8.4554193489\,I-G$ are positive and $\det(8.4554193488\,I-G)<0$, which places the largest eigenvalue, a root of $Q_{13}$ because it exceeds $3$, in the stated interval; it is therefore $\beta$. This completes the proof. \qed

\begin{remark}
The reduction can be compared with direct computation at primes outside $E$. At $p=17,19,23,101,257$ and $1009$ the characteristic polynomial of the Gram matrix at every centre is one of the $23$ model polynomials, and the maximal one occurs only at $x=1$. A floating-point computation of $N_2(p)$ at every prime up to $200$ and at $257$ and $1009$ agrees with \cref{thm:main}.
\end{remark}

\begin{remark}
The theorem concerns prime fields. For $\F_q$ with $q$ a prime power the short stabilizers in small characteristic would have to be analysed again. The input space is the full coordinate space on the ball; restricting the output to the ball, or the input to vectors of sum zero, changes the extremal problem.
\end{remark}

\section{Computations}\label{sec:computations}

The program \texttt{certificate.py} in the repository below uses only the Python standard library, and every sign decision in it is made with integers or with \texttt{fractions.Fraction}. It enumerates $B_0,\ldots,B_6$, computes $\Delta$ on all pairs, and checks that singular pairs commute. It finds $E$ from common fixed points at every point of $X_p$ for $p\le101$, and the scalar-reduction primes. It forms the $22$ commutation classes and the $23$ models, and checks the displayed Gram matrix of $C_*$ up to a permutation of coordinates. Characteristic polynomials are computed by the Faddeev--LeVerrier recursion over the rationals. The Sylvester certificates for the other models and the brackets for $\alpha$, $\sqrt\alpha$ and $\beta$ are checked by exact minors and determinants. The program then carries out the computations at the primes of $E$ described in \cref{sec:proof} and the comparison in the first remark. It finishes in under a minute on a laptop and writes its results to \texttt{certificate\_results.json}. The program \texttt{float\_check.py} is the floating-point check, using NumPy~\cite{NumPy}.

\section*{Acknowledgments and funding}
The research presented in this paper was supported by the European Research Council (ERC) under the European Union's Horizon Europe research and innovation programme (grant agreement No.~101041711), by the Simons Foundation as part of the Collaboration on the Mathematical and Scientific Foundations of Deep Learning, by Heights Labs, by Convex Nexus Capital, by the Israel Science Foundation (grant number 2258/19), and by the Israel Science Foundation (ISF Grant 4101/25).

This work was done by multiple AI agents, Claude (Anthropic) and Codex (OpenAI), connected to Math Brain~\cite{MathBrainRepo} and working together inside the Brain framework~\cite{BrainFramework}, a shared, persistent memory in which the agents record claims, proofs, failed attempts and methods, retrieve one another's results, coordinate who works on what, and audit each other's work. The method is described in~\cite{MathBrainPaper}.

\section*{Declaration of competing interest}
The author holds equity in Heights Labs and in Convex Nexus Capital and has a management role at Convex Nexus Capital. These entities are included in the funding declaration above.

\section*{Code and data availability}
The \LaTeX\ source, the exact certificate program and its output, and the floating-point check are available at \url{https://github.com/yspennstate/modular-projective-line-radius-two-norm/tree/v1.0}.

\section*{Declaration of generative AI and AI-assisted technologies}
Codex (OpenAI) found the reduction to short stabilizers and the $23$ local models and computed the first exact certificates; a second Codex run checked them with separate code. Claude (Anthropic) wrote the certificate program and the floating-point check published with the paper, independently of that code, and wrote the text. The question is the deterministic one posed by van Handel~\cite{vanHandel} and conjectured by Buck~\cite{Buck} for diffusion operators. The regular norm goes back to Cartwright and Soardi~\cite{CartwrightSoardi}. The common-eigenline criterion is due to Friedland~\cite{Friedland}, in the form given by Florentino~\cite{Florentino}. The author selected the problem, directed the work, and is responsible for the content of the paper and for decisions concerning its dissemination.

\begin{thebibliography}{99}

\bibitem{BarbieriJezernik} M.~Barbieri and U.~Jezernik, \emph{Strong convergence of random representations of free products of finite groups}, arXiv:2602.11905, 2026, \url{https://arxiv.org/abs/2602.11905}.

\bibitem{Buck} M.~W.~Buck, \emph{Expanders and diffusers}, SIAM J. Algebraic Discrete Methods \textbf{7} (1986), 282--304, doi:10.1137/0607032.

\bibitem{CartwrightSoardi} D.~I.~Cartwright and P.~M.~Soardi, \emph{Random walks on free products, quotients and amalgams}, Nagoya Math. J. \textbf{102} (1986), 163--180, doi:10.1017/S0027763000000507.

\bibitem{Florentino} C.~A.~A.~Florentino, \emph{Simultaneous similarity and triangularization of sets of $2$ by $2$ matrices}, Linear Algebra Appl. \textbf{431} (2009), 1652--1674; arXiv:0809.3032.

\bibitem{Friedland} S.~Friedland, \emph{Simultaneous similarity of matrices}, Adv. Math. \textbf{50} (1983), 189--265, doi:10.1016/0001-8708(83)90044-0.

\bibitem{GolubevKamber} K.~Golubev and A.~Kamber, \emph{Cutoff on graphs and the Sarnak--Xue density of eigenvalues}, arXiv:1905.11165.

\bibitem{LaffertyRockmore92} J.~D.~Lafferty and D.~Rockmore, \emph{Fast Fourier analysis for $\mathrm{SL}_2$ over a finite field and related numerical experiments}, Experiment. Math. \textbf{1} (1992), 115--139, doi:10.1080/10586458.1992.10504252.

\bibitem{LaffertyRockmore93} J.~Lafferty and D.~Rockmore, \emph{Numerical investigation of the spectrum for certain families of Cayley graphs}, in: Expanding Graphs, DIMACS Ser. Discrete Math. Theoret. Comput. Sci. \textbf{10}, Amer. Math. Soc., 1993, 63--73, doi:10.1090/dimacs/010/06.

\bibitem{RivinSardari} I.~Rivin and N.~T.~Sardari, \emph{Quantum chaos on random Cayley graphs of $\mathrm{SL}_2[\mathbb{Z}/p\mathbb{Z}]$}, Exp. Math. \textbf{28} (2019), 328--341, doi:10.1080/10586458.2017.1403982.

\bibitem{MathBrainRepo} Y.~Shmalo, \emph{Math Brain: public snapshot, code and tests}, 2026, \url{https://github.com/yspennstate/math-brain-public/tree/5427734deee9be643d6dfbdd5e7d4baeeefe4a4a}.

\bibitem{MathBrainPaper} Y.~Shmalo, \emph{Math Brain: shared, auditable memory for multi-agent mathematical research}, October 2026, doi:10.5281/zenodo.24150837, \url{https://github.com/yspennstate/math-brain-public/tree/5427734deee9be643d6dfbdd5e7d4baeeefe4a4a/papers/math-brain-2026-10-24}.

\bibitem{BrainFramework} Y.~Shmalo, \emph{Brain Framework}, public software repository, 2026, \url{https://github.com/yspennstate/brain-framework-public/tree/46b935d26d5a2893f7790b60466387dd8da47ef5}.

\bibitem{NumPy} C.~R.~Harris, K.~J.~Millman, S.~J.~van der Walt, R.~Gommers, P.~Virtanen, D.~Cournapeau, E.~Wieser, J.~Taylor, S.~Berg, N.~J.~Smith, R.~Kern, M.~Picus, S.~Hoyer, M.~H.~van Kerkwijk, M.~Brett, A.~Haldane, J.~Fern\'andez del R\'io, M.~Wiebe, P.~Peterson, P.~G\'erard-Marchant, K.~Sheppard, T.~Reddy, W.~Weckesser, H.~Abbasi, C.~Gohlke and T.~E.~Oliphant, \emph{Array programming with NumPy}, Nature \textbf{585} (2020), 357--362, doi:10.1038/s41586-020-2649-2.

\bibitem{vanHandel} R.~van Handel, \emph{The strong convergence phenomenon}, arXiv:2507.00346, 2025, \url{https://arxiv.org/abs/2507.00346}.

\end{thebibliography}
\end{document}
